%======================================================================
\section{The marked-pair identity: eliminating the post-flip events}
\label{sec:marked}

The identity of Theorem~\ref{thm:identity} expresses the completion
ratio through the post-flip events $B_1,B_2$ of the \emph{joining}
measure.  In this section we prove a cleaner exact identity that
eliminates the joining measure, the switch tie-break, and the events
$B_1,B_2$ altogether: the ratio $\CC_n(\lambda)/\CC_n(\mu)$ is the
$A$-versus-$B$ \emph{odds of a single marked column pair in a single
uniform square of} $S(\lambda)$.  The proof is self-contained (it does
not use the multigraph of \cite{CGW08}), covers $\alpha=\beta$ without
any modification --- resolving the bookkeeping discrepancy of
Remark~\ref{rem:albe} --- and yields the lower bound
$\CC_n(\lambda)/\CC_n(\mu)>\tfrac12$ for free.

Throughout, squares are normalised to have first row $\mathrm{id}$, so
the column map is $\omega=\sigma$ ($1\circ\omega(j)=2\circ j$ with the
first row the identity gives $\omega(j)=\sigma(j)$).  For the proofs it
is convenient to work with \emph{unnormalised} squares (equivalently,
symbol-relabelling classes): every instance count below then acquires a
common factor $n!$, which cancels from every ratio, and the maps we use
(flip, switch) act on unnormalised squares without any renormalisation
step.  Statuses, cycle types, and column positions are
relabelling-invariant, so we freely state results in the normalised
picture.

\begin{definition}[marked instances]\label{def:marked}
Let $\lambda\vdash n$ contain a part $m=\alpha+\beta$
($\alpha,\beta\ge2$), and let $\mu$ be $\lambda$ with one part $m$
split into $\alpha,\beta$.  The \emph{splitting instance space} is
\[
X \;=\; X(n,\lambda;\alpha,\beta)
\;=\;\bigl\{(L',P)\;:\;L'\in S(\lambda),\;
P=\{j,\sigma^\alpha(j)\}\text{ with } j\text{ in an $m$-cycle of }
\sigma(L')\bigr\},
\]
unordered pairs counted once.  Thus every $L'$ carries exactly
$\lambda_m\, m$ marked pairs when $\alpha\neq\beta$ and
$\lambda_m\,\alpha$ when $\alpha=\beta$.  Since $\alpha,\beta\ge 2$
every marked pair is admissible, and $X=X_A\sqcup X_B$ according to the
$A$/$B$ status of $P$ in $L'$.  The \emph{joining instance space} is
\[
J \;=\; J(n,\mu;\alpha,\beta)
\;=\;\bigl\{(L,P)\;:\;L\in S(\mu),\;P=\{j,j'\},\;
j\text{ in an $\alpha$-cycle},\; j'\text{ in a different
$\beta$-cycle}\bigr\},
\]
again with $J=J_A\sqcup J_B$ by the status of $P$ in $L$.
\end{definition}

The counting arithmetic is immediate: with
$D_\nu$ the number of derangements of cycle type $\nu$ and
$N(\nu)=\CC_n(\nu)(n-2)!$,
\begin{equation}\label{eq:Xarith}
\frac{2|X|}{|J|}
=\frac{2\,D_\lambda N(\lambda)\,\lambda_m m}
       {D_\mu N(\mu)\,\mu_\alpha\mu_\beta\,\alpha\beta}
=\frac{2\CC_n(\lambda)}{\CC_n(\mu)}
\qquad(\alpha\neq\beta),
\end{equation}
using $D_\lambda/D_\mu=\alpha\mu_\alpha\,\beta\mu_\beta/(m\lambda_m)$;
for $\alpha=\beta$ one has
$D_\lambda/D_\mu=\alpha^2\mu_\alpha(\mu_\alpha-1)/(2m\lambda_m)$,
$|J|=|S(\mu)|\binom{\mu_\alpha}{2}\alpha^2$ and
$|X|=|S(\lambda)|\lambda_m\alpha$, and \eqref{eq:Xarith} holds
verbatim.  The content of the identity is therefore the following
count.

\begin{lemma}[cycle trades preserve pair status]\label{lem:cycletrade}
Let $P=\{j,j'\}$ be an admissible column pair of a Latin square $L$ and
let $C$ be one column cycle of $P$.  Trading $C$ (swapping the entries
of columns $j,j'$ in the rows of $C$) inverts the pair permutation
$\pi_P$ on $C$ and fixes it elsewhere; the partition of rows into
column cycles of $P$ is unchanged.  In particular the $A$/$B$ status of
$P$ itself is invariant under the CGW flip at $P$, which trades the
full column cycle through row~$2$.
\end{lemma}

\begin{proof}
Write $\pi=\pi_P$, so $\pi(r)$ is the row of column $j'$ holding the
symbol $L(r,j)$; $C$ is $\pi$-closed and the trade swaps
$L(r,j)\leftrightarrow L(r,j')$ for $r\in C$.  Fix $r\in C$ and let
$\pi'$ be the new pair permutation.  The new $L(r,j)$ is the old
$L(r,j')$; the row of the new column $j'$ holding that symbol is found
by solving (new $L(s,j')$) $=$ (old $L(r,j')$): for $s\in C$ the new
$L(s,j')$ is the old $L(s,j)$, and old $L(s,j)=$ old $L(r,j')$ exactly
when $\pi(s)=r$.  Hence $\pi'=\pi^{-1}$ on $C$; off $C$ nothing
changes.  A cycle inverted in place has the same support, so the
partition into cycles is unchanged.
\end{proof}

\begin{lemma}[switch toggle]\label{lem:switchtoggle}
Let $(L,P)\in J$ and let $Q$ be the row cycle through one column of
$P$ (the two row cycles are distinct).  Trading $Q$ (the CGW switch)
changes $\pi_P$ by precomposition with the transposition of rows
$1,2$: the cycles of $\pi_P$ through rows $1,2$ are split if rows
$1,2$ are in one cycle, and joined if in two.  Hence the switch
toggles the status of $P$, preserves $S(\mu)$ and the instance data
(it inverts one $\sigma$-cycle in place), and is an involution:
the trade preserves the traded cycle's symbol set (the two rows carry
the same symbols on its columns), so the min-symbol rule re-selects the
same cycle.  Consequently $|J_A|=|J_B|$, exactly.
\end{lemma}

\begin{proof}
Say $Q$ is the row cycle through $j$; its columns include $j$ but not
$j'$.  The trade swaps the row-1 and row-2 entries in the columns of
$Q$; for $\pi_P$ only column $j$ matters, where the entries of rows
$1,2$ are exchanged.  Thus $\pi_P'=\pi_P\circ\tau_{12}$, and
multiplying by a transposition splits or joins the cycles through its
two points.  Rows $1,2$ lie in one cycle of $\pi_P$ exactly when $P$
is a $B$-pair.  The trade swaps the two rows' entries along the
columns of $Q$, which inverts the corresponding $\sigma$-cycle in
place: the cycle type and the column sets of all row cycles are
preserved, so $(L,P)\mapsto(\mathrm{sw}(L),P)$ maps $J$ to $J$.
\end{proof}

\begin{theorem}[marked-pair identity]\label{thm:marked}
For every $n$, $\lambda$, and split $m=\alpha+\beta$ with
$\alpha,\beta\ge2$ (including $\alpha=\beta$),
\[
|J_A|\;=\;|J_B|\;=\;|X_A|,
\qquad\text{hence}\qquad
\boxed{\;\frac{\CC_n(\lambda)}{\CC_n(\mu)}\;=\;\frac{|X|}{2|X_A|},
\qquad
2\,\frac{\CC_n(\lambda)}{\CC_n(\mu)}-1\;=\;\frac{|X_B|}{|X_A|}.\;}
\]
In particular $\CC_n(\lambda)>\tfrac12\,\CC_n(\mu)$ unconditionally,
and Hypothesis~\ref{hyp:eq} ($\mathrm{EQ}(\delta)$) is equivalent to
\[
\Bigl|\;\Pr\bigl(P\text{ is a $B$-pair}\bigr)-\tfrac12\;\Bigr|
\;\le\;\tfrac{\delta}{4}(1+o(1)),
\]
for $(L',P)$ uniform on $X$: the law of a single bit --- whether rows
$1,2$ lie in one column cycle of the marked pair
$\{j,\sigma^\alpha(j)\}$ --- of a single uniform square of
$S(\lambda)$.
\end{theorem}

\begin{proof}
By \eqref{eq:Xarith} it suffices to prove $|J|=2|X_A|$.
The flip at an $A$-pair $P$ whose columns lie in one row cycle at
$\sigma$-distances $(\alpha,\beta)$ splits that cycle into an
$\alpha$- and a $\beta$-cycle and maps $S(\lambda)\to S(\mu)$; by
Lemma~\ref{lem:cycletrade} the pair $P$ is still an $A$-pair
afterwards, so $x\mapsto\mathrm{unflip}(x)$ maps $X_A$ into $J_A$.
The flip at a fixed pair is an involution (the traded cycle again
contains row 2, with the same row support), so
$\mathrm{join}=\mathrm{flip}$ inverts it: $X_A\leftrightarrow J_A$ is
a bijection.  Composing with the switch of
Lemma~\ref{lem:switchtoggle} gives a bijection
$X_A\leftrightarrow J_B$, namely
$x\mapsto \mathrm{sw}(\mathrm{unflip}(x))$ with inverse
$(L,P)\mapsto \mathrm{flip}(\mathrm{sw}(L))$ --- the map
``switch-then-flip'', which is exactly the CGW joining step at a
$B$-pair.  Hence $|J|=|J_A|+|J_B|=2|X_A|$.
\end{proof}

\begin{remark}[what became of the events $B_1,B_2$]
Theorem~\ref{thm:marked} recovers Theorem~\ref{thm:identity} and
sharpens it: the joining measure, the min-symbol averaging, and the
post-flip events have disappeared.  The mechanism is that
\emph{every} joining edge lands in $X_A$ (Lemma~\ref{lem:cycletrade}),
each marked-$A$ instance receiving exactly two (one from an $A$-join,
one from its switched partner), while $X_B$ is never reached: the
$B_1,B_2$ bookkeeping of \cite{CGW08} was the joining-side accounting
of $|X_B|$.  On $X_B$ itself the only known measure-preserving
involution is the cross-switch of \cite[Def.~3.9]{CGW08}.  By
\cite[Lem.~3.10]{CGW08} it toggles the status of every pair
$\{\omega^{k}(j),\omega^{k'}(j')\}$ with $1\le k<\alpha$,
$1\le k'<\beta$ --- a family containing \emph{both} shifted pairs, at
$(k,k')=(1,1)$ and $(\alpha-1,\beta-1)$ --- while (empirically: all
$2880$ marked-$B$ instances at $n=5$) fixing the marked status
itself.  This yields
$\Pr(B_1\mid\text{marked }B)=\Pr(B_2\mid\text{marked }B)=\tfrac12$
exactly, explaining the extraordinary balance of the post-flip events
observed in \S\ref{sec:identity}.
\end{remark}

\begin{remark}[the case $\alpha=\beta$, and Remark~\ref{rem:albe}]
Nothing in the proof distinguishes $\alpha=\beta$: marked pairs
$\{j,\sigma^\alpha(j)\}$ in a $2\alpha$-cycle are counted once
($\lambda_m\alpha$ per square), and \eqref{eq:Xarith} holds with the
$\alpha=\beta$ conventions.  Exhaustive verification at $n=6$:
$\mu=(3,3)$ gives $|X_B|/|X_A|=622080/829440=3/4$, exactly the value
$2\CC_6((6))/\CC_6((3,3))-1$ that the joining-side bookkeeping of
Remark~\ref{rem:albe} failed to reproduce.  The $\alpha=\beta\ge3$
anomaly is thus resolved: the identity was never about the joining
events.
\end{remark}

\paragraph{Verification.}  By complete enumeration
(\texttt{pf\_bridge.py}, \texttt{xside.c}): at $n=5$,
$\mu=(2,3)$: $|J_A|=|J_B|=|X_A|=1440$, $|X_B|=2880$, ratio $2$
$=\bar q$ exactly.  At $n=6$: $(2,4)\to(6)$:
$1520640/1520640/1520640$, $|X_B|=1382400$, ratio $10/11$;
$(2,2,2)\to(2,4)$: ratio $4/7$; $(3,3)\to(6)$: ratio $3/4$ --- all
equal to $2\CC_6(\lambda)/\CC_6(\mu)-1$ exactly.  At $n=7$, fixing one
representative $\sigma_0$ per type (conjugation symmetry) and
enumerating all $\approx6.6\times10^6$ completions per type
(\texttt{xside2.c}), the $X$-side odds reproduce
$2\CC_7(\lambda)/\CC_7(\mu)-1$ exactly for all four admissible merges:
$(2,5)\to(7)$: $22848000/23116800=85/86$;
$(3,4)\to(7)$: $23063040/22901760=143/142$ (note the $B$-excess);
$(2,2,3)\to(2,5)$: $107/108$; $(2,2,3)\to(3,4)$: $35/36$.
The bijections themselves
(unflip status, roundtrips, hit multiplicities, orphan structure) were
verified instance-by-instance at $n=5,6$.
